\documentclass[12pt]{article}
\usepackage[T1]{fontenc}
\usepackage[utf8]{inputenc}
\usepackage{lmodern}
\usepackage{textcomp}
\usepackage{enumitem}
\usepackage{geometry}
\usepackage{graphicx}
\usepackage{amsmath, amssymb}
\geometry{margin=1in}
\begin{document}
\begin{center}
Physics - Undergraduate Level Assessment 4 \
Total Marks: 40 \
Answer all questions.
\end{center}

\section*{Multiple Choice (10 marks)}
\begin{enumerate}[label=\arabic*.]
\item Electromagnetic radiation emitted from a nucleus is most likely to be in the form of
\begin{enumerate}[label=(\alph*)]
    \item gamma rays
    \item microwaves
    \item ultraviolet radiation
    \item visible light
\end{enumerate}
\item An atom has filled n = 1 and n = 2 levels. How many electrons does the atom have?
\begin{enumerate}[label=(\alph*)]
    \item 2
    \item 4
    \item 6
    \item 10
\end{enumerate}
\item A resistor in a circuit dissipates energy at a rate of 1 W. If the voltage across the resistor is doubled, what will be the new rate of energy dissipation?
\begin{enumerate}[label=(\alph*)]
    \item 0.25 W
    \item 0.5 W
    \item 1 W
    \item 4 W
\end{enumerate}
\item Which of the following statements about bosons and/or fermions is true?
\begin{enumerate}[label=(\alph*)]
    \item Bosons have symmetric wave functions and obey the Pauli exclusion principle.
    \item Bosons have antisymmetric wave functions and do not obey the Pauli exclusion principle.
    \item Fermions have symmetric wave functions and obey the Pauli exclusion principle.
    \item Fermions have antisymmetric wave functions and obey the Pauli exclusion principle.
\end{enumerate}
\item The rest mass of a particle with total energy 5.0 GeV and momentum 4.9 GeV/c is approximately
\begin{enumerate}[label=(\alph*)]
    \item 0.1 GeV/$c^2$
    \item 0.2 GeV/$c^2$
    \item 0.5 GeV/$c^2$
    \item 1.0 GeV/$c^2$
\end{enumerate}
\end{enumerate}

\section*{True / False (10 marks)}
\textit{Answer True or False}
\begin{enumerate}[label=\arabic*.]
\setcounter{enumi}{5}
\item True or False: A helium-neon laser is the most versatile option for spectroscopy in the visible range.
\item True or False: The negative muon, mu^-, has properties most similar to the electron, including charge and lepton classification.
\item True or False: The angular magnification of the refracting telescope is 4 when the eye-piece lens has a focal length of 20 cm.
\item True or False: Boron cannot be used as a dopant in germanium to create an n-type semiconductor.
\item True or False: A gas laser operates by utilizing transitions between the energy levels of free atoms in a gaseous medium.
\end{enumerate}

\section*{Long Form Response (20 marks)}
\textit{Answer each question with a clear and complete explanation.}
\begin{enumerate}[label=\arabic*.]
\setcounter{enumi}{10}
\item Discuss the behavior of unpolarized light as it passes through two ideal linear polarizers set at a 45-degree angle to each other, and analyze the resulting intensity of the transmitted light in relation to the incident intensity.
\item Discuss the distribution of kinetic energy in a uniform solid disk rolling down an inclined plane without slipping, and analyze the fraction that is rotational kinetic energy after a given time.
\end{enumerate}

\vfill
\noindent\textit{End of Paper}
\end{document}